\id{ҒТАМР 50.47.02}{}

{\bfseries ЦИФРЛЫҚ ЕГІЗ АРХИТЕКТУРАСЫНДАҒЫ ДЕРЕКТЕР АҒЫНЫН ОҢТАЙЛАНДЫРУ:
IIOT ХАТТАМАЛАРЫНЫҢ ӨНІМДІЛІГІ МЕН ҚАУІПСІЗДІГІН КЕШЕНДІ ТАЛДАУ}

{\bfseries Г.А.Амирханова}\fig{i2/image1}[]{\bfseries ,
Е.Фарузқызы}\fig{i2/image1}[]{\bfseries \envelope ,
Б.С.Амирханов}\fig{i2/image1}[]{\bfseries ,
А.А.
Раева}\fig{i2/image1}[]{\bfseries ,
Д.С.
Жайсанова}\fig{i2/image1}[]

\emph{\tsp{1}әл-Фараби атындағы Қазақ ұлттық университеті,
Алматы, Қазақстан}

\envelope \emph{Корреспондент-автор:
elnura.dauletkhan@gmail.com}

Өндіріс циклін басқаруда цифрлық егіз технологиясын пайдалану
компаниялар үшін тығыз желілерде жоғары жылдамдықты деректерді берумен
байланысты қиындықтарға және қауіпсіздіктің жоғары деңгейін қамтамасыз
ету қажеттілігіне байланысты өте маңызды. Бұл зерттеуде BME280
сенсорлары мен ESP32-S3 микроконтроллерін қамтитын архитектурадағы үш
өнеркәсіптік интернет заттарының (IIoT) хаттамаларының (MQTT 5.0,
WebSocket және HTTP/2) операциялық тиімділігі мен қауіпсіздік деңгейін
терең бағалау ұсынылған. Зерттеу сонымен қатар деректерді беру үшін
екілік хаттама буферлік кодтау форматын (Protobuf) пайдаланудың
артықшылығы туралы объективті дәлелдер келтіреді. Бұл зерттеуде
қолданылған әдіснама деректерді өткізу қабілетін, екі жақты сапар
уақытын (RTT) және ресурстар жүктемесін кешенді бағалауды қамтыды.
Қауіпсіздікті жақсарту үшін Transport Layer Security (TLS) 1.3
хаттамасының өңдеу уақытына және TLS өңдеу үшін аппараттық қауіпсіздік
модульдерін (HSM) пайдалануға әсерін тексеру үшін эмпирикалық талдау
жүргізілді. Сонымен қатар, желілік трафикті басқару үшін шеткі есептеу
деңгейіндегі бейімделгіш сүзгілеу алгоритмдері пайдаланылды. Тәжірибелік
нәтижелер MQTT 5.0 қазіргі уақытта басқа екі хаттамаға қарағанда төмен
кідіріспен қамтамасыз ететінін, ал WebSocket веб-қосымшалардың
тұрақтылығын қамтамасыз ететінін көрсетеді. Бұл зерттеудің негізгі
ғылыми жаңалығы - Protobuf екілік деректер инкапсуляциясын пайдаланатын
көп қабатты IIoT архитектурасын ұсыну. Дәстүрлі JSON-мен салыстырғанда,
бұл деректер алмасу әдісі желі жүктемесін 67,06\%-ға немесе одан да
көпке азайтып, желінің кибершабуылдарға төзімділігін жақсарта алады. Бұл
зерттеудің нәтижелері интеллектуалды және нақты уақыт режимінде жүйені
бақылауға қабілетті кәсіпорынның сандық модельдерін жобалау үшін
стратегиялық негіз болып табылады.

{\bfseries Түйін сөздер:} цифрлық егіз, IIoT хаттамалары, MQTT
5.0,WebSocket, Protobuf, TLS 1.3, Edge Computing, индустрия 4.0.

{\bfseries ОПТИМИЗАЦИЯ ПОТОКА ДАННЫХ В АРХИТЕКТУРЕ ЦИФРОВОГО ДВОЙНИКА:
КОМПЛЕКСНЫЙ АНАЛИЗ ПРОИЗВОДИТЕЛЬНОСТИ И БЕЗОПАСНОСТИ ПРОТОКОЛОВ IIOT}

{\bfseries Г.А.Амирханова, Е.Фарузкызы\envelope , Б.С.Амирханов,~
А.А Раева, Д.С Жайсанова}

\emph{Казахский национальный университет имени аль-Фараби, Алматы,
Казахстан,}

\emph{e-mail:
elnura.dauletkhan@gmail.com}

Использование технологии цифровых двойников в управлении
производственным циклом имеет жизненно важное значение для компаний
из-за проблем, связанных с передачей данных с высокой скоростью в
плотные сети, и необходимости обеспечения высокого уровня безопасности.
В данном исследовании представлена углубленная
оценка операционной эффективности и уровня безопасности трех протоколов
промышленного интернета вещей (IIoT) (MQTT 5.0, WebSocket и HTTP/2) в
архитектуре, включающей датчики BME280 и микроконтроллер ESP32-S3.
Исследование также предоставляет объективные доказательства
предпочтительного использования двоичного формата кодирования
протокольных буферов (Protobuf) для передачи данных. Методология,
использованная в данном исследовании, представляла собой полную оценку
пропускной способности данных, времени кругового пути (RTT) и нагрузки
на ресурсы. Для повышения уровня безопасности был проведен эмпирический
анализ для проверки влияния протокола Transport Layer Security (TLS) 1.3
на время обработки и влияния использования аппаратных модулей
безопасности (HSM) для обработки TLS. Кроме того, для управления сетевым
трафиком использовались адаптивные алгоритмы фильтрации на уровне
граничных вычислений. Экспериментальные результаты показывают, что MQTT
5.0 в настоящее время обеспечивает меньшую задержку, чем два других
протокола, а WebSocket обеспечивает стабильность для веб-приложений.
Основная научная инновация этого исследования --- предложение
многоуровневой архитектуры IIoT с использованием бинарной инкапсуляции
данных Protobuf. По сравнению с традиционным JSON, использование этого
метода обмена данными может снизить нагрузку на сеть на 67,06\% и более
и повысить устойчивость сети к кибератакам. Результаты этого
исследования обеспечивают стратегическую основу для проектирования
корпоративных цифровых моделей, которые являются «умными» и способны
обеспечивать мониторинг систем в реальном времени.

{\bfseries Ключевые слова:} цифровой двойник, протоколы IioT, MQTT 5.0,
WebSocket, Protobuf, TLS 1.3, Edge Computing (граничные вычисления),
индустрия 4.0.

{\bfseries OPTIMIZATION OF DATA FLOW IN DIGITAL TWIN ARCHITECTURE: A
COMPREHENSIVE ANALYSIS OF PERFORMANCE AND SECURITY OF IIOT PROTOCOLS}

{\bfseries G.A.Amirkhanova, Y.Faruzkyzy\envelope , B.S.Amirkhanov,
A.A.Rayeva, D.S.Zhaissanova}

\emph{Al-Farabi Kazakh National University, Almaty, Kazakhstan,}

\emph{e-mail:
elnura.dauletkhan@gmail.com}

It is vital for companies to use Digital Twin technology in the
management of their production cycle due to challenges of data being
transferred at high rates into dense networks and the need for high
levels of security. This research presents an in-depth evaluation of the
operational efficiency and level of security of three industrial
internet of things (IIoT) protocols (MQTT 5.0, WebSocket, and HTTP/2)
within an architecture incorporating BME280 sensors and ESP32-S3
microcontroller. The study also provides objective evidence for the
preferred use of the protocol buffers (Protobuf) binary encoding format
for transmitting data.The methodology utilized in this research was a
complete assessment of data throughput, round-trip time (RTT), and
resource load. In order to increase the level of security, an empirical
analysis was performed to test the effect of Transport Layer Security
(TLS) 1.3 with regard to processing time and the impact of the use of
Hardware Security Modules (HSM) to process TLS. Additionally, adaptive
filtering algorithms at the Edge Computing level were employed to
control network traffic.The experimental results demonstrate that MQTT
5.0 currently produces less latency than the other two protocols, while
WebSocket provides stability for web applications. The primary
scientific innovation of this study is the proposition of a multi-level
IIoT architecture utilizing Protobuf binary data encapsulation. When
compared to traditional JSON, the use of this method of data exchange
can reduce network load by 67.06 percent or more and increase network
resilience against cyberattacks. The results of this study provide a
strategic foundation for the design of enterprise digital models that
are "smart" and capable of providing real-time monitoring of systems.

{\bfseries Keywords:} digital twin, IIoT protocols, MQTT 5.0, WebSocket,
Protobuf, TLS 1.3, Edge Computing, industry 4.0.

{\bfseries Кіріспе.} Қазіргі уақытта әлемдік бизнес секторы «Индустрия 4.0»
парадигмасынан туындаған елеулі техникалық реформалық өзгерістерді
бастан кешіруде. Бұл трансформацияның негізгі компоненттерінің бірі -
материалдық активтердің динамикалық виртуалды көрінісін жасайтын сандық
егіздер. Сандық егіздер тек визуалды көрініс үшін ғана емес, сонымен
қатар активтердің болашақ болжамдарын әзірлеу, активпен не болғанын тез
тіркеу және технологияның үздіксіз жұмыс істеу тәсілін жақсартуға
көмектесу үшін толық интеллектуалды негіз болып табылады.

Дегенмен, сандық егіздердің функционалдығы физикалық объект пен оның
виртуалды көрінісі арасындағы ақпараттың жылдам және дәл берілуіне
айтарлықтай байланысты. Қазіргі заманғы зауыттарда мыңдаған сенсорлық
түйіндер айтарлықтай көлемде деректер жасайды, бұл желілік
инфрақұрылымға жүктеме түсіреді, кідіріс уақытын арттырады және деректер
қауіпсіздігін төмендетеді. Сондықтан, тиісті байланыс хаттамаларын
таңдау және деректерді сериялау форматтарын оңтайландыру Заттар
өнеркәсіптік интернеті (IIoT) саласында зерттеу және қолданба тұрғысынан
маңызды бола түсуде.

\emph{Зерттеудің мақсаты мен жаңалығы.} Бұл зерттеудің мақсаты -
ресурстармен шектелген сандық егіздердің өнеркәсіптік интернет заттары
(IIoT) құрылғылары үшін оңтайлы деректер ағынының архитектурасын
анықтау. Жобада IIoT құрылғыларынан телеметрияны қабылдау үшін
аппараттық құрал ретінде ESP32-S3 микроконтроллері және BME280 сенсоры
қолданылады. Бұл зерттеу құрылымдық деректерді сериялауға арналған
екілік формат Protobuf-тың дәстүрлі түрде қолданылатын JSON форматымен
салыстырғанда артықшылықтарын көрсетеді, бұл дәстүрлі телеметрия
деректерін жіберу әдістерімен салыстырғанда пайдалы жүктеме көлемін
айтарлықтай азайтады.

\fig{i2/image17}[]

{\bfseries 1-сурет. BME280 сандық сенсоры және ESP32-S3 микроконтроллері}

Зерттеу үшін таңдалған аппараттық жүйе (1-сурет) нақты уақыт режиміндегі
деректерді жинауға және өңдеуге арналған. Қоршаған орта температурасын,
ылғалдылығын және атмосфералық қысымды жоғары дәлдікпен өлшеу үшін
BME280 сенсоры қолданылады. ESP32-S3 микроконтроллері жиналған
деректерді MQTT 5.0 немесе WebSocket арқылы сандық егіздердің виртуалды
моделіне жібермес бұрын TLS 1.3 хаттамасын пайдаланып шифрлайды.

Осы мақалада ұсынылған қосымша ғылыми инновация ретінде шеткі есептеу
деңгейінде шекті бейімделгіш сүзгілеу әдісі әзірленді. Protobuf
форматымен біріктірілген бұл әдіс өлшенген уақыт параметрлеріндегі
физикалық маңызды ауытқуларды (шу деңгейінен тыс ауытқулар) ғана жіберу
арқылы құрылғылар арасындағы желілік трафикті оңтайландырудың жаңа
тәсілін ұсынады.

\emph{Эксперименттік аппараттық кешен.} Әдеттегі өндірістік орталардың
дәлірек көрінісін алу үшін біз осы зерттеуге әртүрлі аппараттық
қабаттарды біріктірдік. Біздің жүйеміз қоршаған орта температурасын,
ылғалдылығын және атмосфералық қысымды жоғары дәлдікпен және 20 биттік
ажыратымдылықпен өлшейтін BME280 сенсорлық модульдеріне негізделген.
Шикі деректерді өңдеу және жүйелі басқару үшін Ryzen 5 процессоры мен
Ubuntu 22.04 LTS операциялық жүйесіне негізделген жоғары өнімді сервер
инфрақұрылымын пайдаланумен қатар, біз бұл серверді деректер ағынын
оңтайландыру және шеткі есептеу деңгейінде деректерді алдын ала сүзу
үшін Mosquitto брокері ретінде де пайдаландық.

\emph{Салыстырмалы хаттамалар талдау.} Үш негізгі ақпарат алмасу әдісін
терең зерттеу жүргізілді, оны кейіннен әртүрлі сандық егіз пайдалану
жағдайлары мен қолданбаларына экстраполяциялауға болады. Деректер
пакетінің құрылымында қолданылатын жеңіл тақырыптардың тиімділігін және
салалық стандартты MQTT 5.0 хаттамасының қызмет көрсету сапасын (QoS)
бағалау үшін бірнеше талдау жүргізілді. Талдау сонымен қатар WebSocket
хаттамасының нақты уақыт режиміндегі 3D визуализация интерфейстері мен
басқару тақталарын әзірлеуге көмектесетін толық дуплексті байланысты
қамтамасыз ету әлеуетін көрсетті. Сонымен қатар, REST API арқылы бұлттық
сақтау орнына жіберілетін үлкен деректер жиынтығын (көптеген деректерді
жинау) мерзімді синхрондаумен байланысты мультиплекстеу мүмкіндіктері
мен ресурстардың тиімділігіне аналитикалық бағалау жүргізілді.

{\bfseries Әдебиеттерге шолу.} Тізімнің {[}1{]} жұмысында өндірісте цифрлық
егіздерді әзірлеу мен пайдаланудағы қазіргі үрдістер мен заманауи
жағдайдың егжей-тегжейлі сипаттамасы берілген. Интернет арқылы сенімді
деректерді беру үшін TCP басым әдіс болып қала береді {[}2{]}, ал
WebSocket нақты уақыт режиміндегі веб-интерфейстер үшін барған сайын
танымал бола бастады {[}3{]}.

MQTT 5.0 нұсқасы ресурстары шектеулі құрылғыларды қосу үшін стандартты
таңдау болып қала береді {[}4{]} және Мишра мен Кертестің зерттеулеріне
сәйкес, басқа хаттамаларға қарағанда энергияны тиімдірек пайдаланады
{[}5{]}. Қауіпсіздік тұрғысынан TLS 1.3 қазір деректерді қорғаудың жаңа
стандарты болып табылады, ескірген TLS-тің бұрынғы нұсқаларын ауыстырады
{[}6{]}.

Заттардың өнеркәсіптік интернетінде (IIoT) әртүрлі деректерді беру
әдістерінің кідірісін зерттеген Сильва және т.б. зерттеуі MQTT HTTP-ге
қарағанда жақсы жауап беру уақытына ие екенін анықтады {[}7{]}.
Дегенмен, Лакшминараян атап өткендей, MQTT хаттамасының қызмет
көрсетуден бас тарту (DoS) шабуылдарына осалдығы мұндай шабуылдардан
қорғану үшін қосымша шифрлау әдістерін қажет етеді {[}8{]}. Сандық
егіздер арасындағы нақты уақыт режиміндегі синхрондаудың мысалдары Чо
мен Ноның {[}9{]} еңбегінде келтірілген.

Жақында іске қосылған maketwin платформасы сандық егіз архитектураның
икемділігін қамтамасыз етеді {[}10{]} және Wi-Fi арнасының (CSI) күй
деректері мен машиналық оқыту технологиялары арқылы нысандарды
бақылаудың жаңа әдістерін енгізеді {[}11{]}. Gao командасы цифрлық
егіздердің кеңсе жиһаздарын басқаруға қаншалықты көмектесетінін
көрсеткенімен {[}12{]}, Скворц WebSockets дуплексті арналар арқылы
деректерді қаншалықты жылдам бөлісе алатынын бағалады {[}13{]}. Михович
әртүрлі хаттамалардың тиімділік коэффициентін бағалау әдісін ұсынады
{[}14{]}; жақында бұл IoT құрылғыларын сайтаралық сценарийлерден (XSS)
және SQL инъекциялық шабуылдардан қалай қорғауға болатынына ғылыми назар
аударуды тудырды {[}15{]}.

OPC UA (open Platform Communications Unified Architecture) ХАТТАМАСЫ
"Индустрия 4.0" компанияларында ақпараттық модельдеу үшін стандарт
ретінде кеңінен қолданылады. OPC UA-ның MQTT-тен артықшылығы-оның сандық
егіздегі ақпараттың семантикалық маңызды сипаттамаларын және оның басқа
сандық егіздермен байланысын қамтамасыз ету мүмкіндігі. Бұл зерттеу OPC
UA және Protobuf интеграциясы арқылы күрделі өнеркәсіптік жүйелердің
цифрлық егіздерін масштабтау мүмкіндігін қарастырады.

Деректерді берудің ең жақсы әдісін таңдау өткізу қабілеттілігін
арттыруға айтарлықтай әсер етеді. Өткізу қабілеті төмен өнеркәсіптік
Заттар интернеті (IIoT) желілерінде Google протоколдық буферлері
(Protobuf) сияқты екілік әдістер JSON сияқты дәстүрлі мәтіндік деректер
форматтарымен салыстырғанда тасымалданатын деректердің көлемін азайтады.
Чо Мен Но MQTT хабар алмасу протоколын қолдана отырып, нақты уақыттағы
синхрондау зауытының цифрлық егізін жасау мүмкін екенін көрсетті.

{\bfseries Материалдар мен әдістер.} Зертханалық қондырғы деректерді беру
үшін digital twin (TSE) архитектурасының тиімділігін арттыру және
арттыру мақсатында жасалған. Деректерді беру/цифрлық егізді құру
процесінің барлық бөліктері қамтылған. Жүйенің жалпы мақсаты-заманауи
TLS v.1.3 шифрлау протоколының көмегімен өндірістік деректерді жылдам
және қауіпсіз алу.

\fig{i2/image18}[]

{\bfseries 2-сурет. Цифрлық егіз жүйесінің деңгейлі архитектурасы}

Бұл зерттеу жобасының архитектурасында төрт архитектуралық деңгей бар.
Барлық төрт деңгей өз функцияларын орындау арқылы деректердің тұтастығы
мен қауіпсіздігін сақтауға көмектеседі.

Бірінші деңгей-сенсорлық деңгей; екіншісі - жетілдірілген есептеу
(edge); үшіншісі - қолданбалы бағдарламалау интерфейсі (API) және осы
зерттеуде қолданылатын бұлтты шешімдер. Сенсор деңгейінде қоршаған
ортаның келесі параметрлерін оқитын аппараттық модуль бар: температура,
ылғалдылық және қысым; және 20 биттік ажыратымдылықпен осы өлшенген
параметрлердің сәйкес сандық көрінісін қамтамасыз етеді. Сенсорлық
қабаттың негізгі элементі-Bosch Sensortec (Германия) шығарған BME280
сандық сенсорлық модулі.

Перифериялық есептеу деңгейі негізінен 240 МГц жиілікте жұмыс істейтін
екі ядросы және бір Xtensa® LX7 бар ESP32 -- S3 микроконтроллерінен
тұрады. Шекаралық есептеу деңгейінде деректерді сериялау протокол
буферлерінің көмегімен өңделеді (Protobuf форматы), ал деректерді
шифрлау TLS 1.3 стандартына сәйкес жүзеге асырылады. AES-256 және RSA
аппараттық үдеткіштерін пайдаланып деректерді шифрлау процесі TLS 1.3
стандартында 0-RTT (Zero Round Trip time Recovernment) шифрлау
протоколдарын пайдалану арқылы шифрланған деректердің жалпы уақыт
кідірісін (кідірісін) азайтады.0-RTT шифрлауды пайдалану бастапқы
құрылғы өшірілген уақыт ішінде "байланысты растау" процесін
орындамай-ақ, олар алынған құрылғыға деректерді лезде жіберуді
қамтамасыз етеді.

Деректерді беру әдісі Wi-Fi желілері арқылы MQTT 5.0 протоколы арқылы
жүзеге асырылады және MQTT хабарламалары Mosquitto дерекқоры арқылы және
оған кіретін Edge шлюзінде орналасқан Mosquitto брокерімен басқарылады.
Соңғы түйін немесе қосымша-бұл Node JS Түйінінің веб-қосымшасы және
WebSocket арқылы MQTT арқылы нақты уақыт режимінде деректерді
визуализациялауға арналған 3D веб-бақылау тақтасы.

{\bfseries 1-кесте. Экспериментте қолданылған аппараттық құралдардың толық
сипаттамасы}

%% \begin{longtable}[]{@{}
%%   >{\raggedright\arraybackslash}p{(\linewidth - 4\tabcolsep) * \real{0.1652}}
%%   >{\raggedright\arraybackslash}p{(\linewidth - 4\tabcolsep) * \real{0.1805}}
%%   >{\raggedright\arraybackslash}p{(\linewidth - 4\tabcolsep) * \real{0.6543}}@{}}
%% \toprule\noalign{}
%% \begin{minipage}[b]{\linewidth}\centering
%% {\bfseries Құрылғы}
%% \end{minipage} & \begin{minipage}[b]{\linewidth}\centering
%% {\bfseries Моделі}
%% \end{minipage} & \begin{minipage}[b]{\linewidth}\centering
%% {\bfseries Техникалық мүмкіндіктері мен зерттеудегі рөлі}
%% \end{minipage} \\
%% \midrule\noalign{}
%% \endhead
%% \bottomrule\noalign{}
%% \endlastfoot
%% {\bfseries Сенсор} & BME280 & Температураны (±0.5°C), ылғалдылықты (±3\%)
%% және қысымды (±1 hPa) өлшеуге арналған кешенді модуль. \\
%% {\bfseries Контроллер} & ESP32-S3 & Деректерді жинауға, сериялауға және
%% аппараттық деңгейде шифрлауға жауапты негізгі есептеуіш блок. \\
%% {\bfseries Шеткі шлюз} & Ryzen 5 3550H / Ubuntu & 16 ГБ жедел жады бар
%% сервер. Ол MQTT брокері ретінде қызмет етеді және деректерді алдын ала
%% терең өңдеуді орындайды. \\
%% {\bfseries Роутер} & Mi Wi-Fi Router 4A & Жергілікті желі ішіндегі байланыс
%% тұрақтылығын және деректер тасымалдау жылдамдығын қамтамасыз етеді. \\
%% \end{longtable}

\fig{i2/image19}[]

{\bfseries 3-сурет. BME280 сенсоры мен ESP32-S3 контроллерінің қосылу
схемасы}

Эксперименттің аппараттық параметрлері төмендегі суретте көрсетілген.
ESP32-S3 BME280-ге I2C интерфейсін (SDA және SCL сызықтары) қолдана
отырып, тақтада қосылды. Жоғарыдағы схема бастапқы деңгейдегі деректерді
алудың және бастапқы деректерді кибершабуылдардан қорғаудың маңыздылығын
көрсетеді. Ryzen 5 процессоры және Ubuntu 22.04 LTS операциялық жүйесі
перифериялық шлюз ретінде пайдаланылады, сонымен қатар жүйенің жалпы
архитектурасының тұрақтылығын қамтамасыз етуге көмектеседі. Осы
параметрді қолдана отырып, жүйе өндіріс жағдайында нақты уақыт режимінде
физикалық параметрлердің стандартты мәндерінен кез-келген ауытқуларды
тез анықтай алады және алынған деректерді виртуалды әлемге қауіпсіз
жібере алады. Бұл процесс деректерді өңдеудің жалпы тиімділігін арттырып
қана қоймайды, сонымен қатар маңызды салалардағы өнеркәсіптік ақпараттық
қауіпсіздіктің барлық қатаң стандарттарына сәйкес келеді.

\emph{Деректерді Protobuf арқылы оңтайландыру.} Google деректерді
өңдеуді тездету және желі арқылы жіберілетін ақпарат көлемін азайту үшін
екілік протоколдық буферлерді (Protobuf) пайдаланатын деректерді
сериялаудың жаңа әдісін жасап шығарды. Protobuf тек өңдеуді жеделдетуге
және есептеу қуатына деген қажеттілікті азайтуға көмектесіп қана
қоймайды; ол деректерге қатысты нақты анықталған типтер жүйесіне ие, ал
JSON форматында мұндай айқын типтер жүйесі жоқ.

Сонымен қатар, Protobuf файлы жоғарыда аталған сенсор деректерінің қалай
ұсынылатынын егжей-тегжейлі сипаттайтын proto деректер құрылымын анықтау
үшін қолданылады. Төменде Protobuf синтаксисін пайдалана отырып, сенсор
хабарламасының құрылымы proto файлында қалай анықталатынына мысал
келтірілген:

Protocol Buffers

syntax = "proto3";

message SensorData \{

float temperature = 1; // 4 байт

float humidity = 2; // 4 байт

float pressure = 3; // 4 байт

int64 timestamp = 4; // 8 байт

\}

Жоғарыдағы деректер типтерінің әрқайсысы үшін қажетті байт санын
пайдалану сол деректерді сақтауға қажетті ең аз орынды (яғни ұзындығын)
анықтауға, сондай-ақ құрылым ішіндегі әрбір өріске бірегей сандық тег
беруге мүмкіндік береді. Мұндай тегтеу әдісі екілік кодпен жазылған
хабарламада деректердің аз мөлшерін сақтауға жағдай жасайды.Пайдалы
жүктеменің (payload) қаншалықты тиімді екенін бағалау үшін келесі
формула қолданылды:

\(S_{reduction} = \frac{{Size}_{JSON}\  - \ {Size}_{Protobuf}}{{Size}_{JSON}} \times 100\%\)
(1)

Астынғы бөлікте көрсетілген диаграммада пакет өлшемінің төмендегені
көрсетілген. Бұл көрсеткіш жүйенің өткізу қабілетіне тікелей әсер етеді:

\fig{i2/image20}[]

{\bfseries 4-сурет. Деректер пакетінің өлшемін салыстыру (JSON vs
Protobuf)}

JSON пакетінде бірдей деректерді (температура, ылғалдылық және қысым)
жіберу үшін орташа өлшемі 80-120 байт қажет. Протокол буферлерінің
(Protobuf) пішімін пайдалану арқылы бұл шамамен 20-30 байтқа дейін
азаяды. Бұл шамамен 67,06\% төмендеуді көрсетеді.

\fig{i2/image21}[]

{\bfseries 5-сурет. Хабарламалар санына байланысты жиынтық трафиктің өсу
динамикасы}

Бұл график уақыт өте келе жіберілетін хабарламалар көлеміне байланысты
жиынтық трафиктің өсуін көрсетеді. Атап айтқанда, нәтижелер JSON (қызыл
сызық) мәтіндік деректер көлемінің экспоненциалды өсуіне әкелетінін және
protobuf (жасыл сызық) протоколы жалпы жүктемені төмендететінін желіде
көрсетеді. Тәжірибе нәтижелері деректерді беру үшін protobuf пішімін
пайдалану трафиктің 67,06\% қысқаруына әкелуі мүмкін деген қорытындыны
қолдайды, бұл өткізу қабілеті шектеулі IIoT жүйелерінің тиімділігін
арттыру және цифрлық егіз технологияларды енгізу үшін өте маңызды.

\fig{i2/image22}[]

{\bfseries 6-сурет. Protobuf қолдану арқылы трафикті үнемдеу тиімділігі}

Талдау көрсеткендей, Protobuf қолданған кезде желілік трафиктің
максималды көлемін 67,06\%- ға төмендетуге болады. Бұл нәтиже өткізу
қабілеті шектеулі өндірістік желілер үшін, сондай-ақ нақты уақыт
режимінде деректерді беруді қажет ететін "сандық егіз" жүйелер үшін өте
маңызды.

Сонымен қатар, екілік форматқа көшу трафикті азайтып қана қоймайды,
сонымен қатар ESP32-S3 микроконтроллерінің қуат тұтынуына оң әсер етеді.
Пайдалы жүктеменің азаюы оның мөлшерінің 3 есе азаюына әкеледі, бұл
Wi-Fi модулінің қуат тұтынуының шамамен 30-40\% төмендейді. Мұндай
энергияны үнемдеуге радионың белсенді тарату режимінде өткізетін уақытын
қысқарту арқылы қол жеткізіледі. Пайдалы жүктеме көлемінің едәуір
азаюына (67,06\%-ға) байланысты ESP32-S3 деректерді тезірек эфирге
жібере алады, бұл оған қысқа уақыт ішінде төмен қуат режиміне ауысуға
мүмкіндік береді. Бұл әдістеме "физикалық объект пен виртуалды модель"
арасындағы деректер ағынын оңтайландыруға және жүйенің жалпы өнімділігін
арттыруға мүмкіндік берді.

\emph{Хаттамаларды конфигурациялау және байланыс модельдері.} Digital
Twin жүйесіндегі ақпарат алмасудың сенімділігін жан-жақты зерттеу үшін
деректер ағынының ерекшеліктеріне бейімделген үш негізгі IIoT хаттамасы
арнайы эксперименттік кезеңде сыналды. Жүйенің негізгі байланыс әдісі
ретінде MQTT 5.0 (TCP арқылы) таңдалды.

Ол Mosquitto брокерін пайдаланып конфигурацияланды және QoS 1 деңгейімен
жабдықталды, бұл қосылымның тұрақтылығын қамтамасыз ету үшін әрбір
пакеттің шлюзге жеткізілуін кепілдендірді. Сонымен қатар, жаңартылған
хаттама мүмкіндіктері деректер тақырыбына метаақпарат қосу арқылы
жүйенің архитектуралық икемділігін айтарлықтай арттырды. Нақты уақыттағы
визуализация деңгейінде RFC 6455 стандартына сәйкес келетін WebSocket
технологиясы пайдаланылды, бұл шеткі есептеу шлюзі мен 3D басқару
тақтасы арасында тұрақты, екі жақты (толық дуплексті) байланыс арнасын
жасады, бұл төмен кідіріспен қамтамасыз етті. Жүйелік файлдарды және
күрделі Protobuf схемаларын мезгіл-мезгіл синхрондау міндеті HTTP/2
хаттамасына жүктелді, бұл мультиплекстеу мүмкіндіктерінің арқасында
желілік ресурстарды тиімді пайдалануға мүмкіндік берді. Байланыс
модельдерін орнату кезінде ESP32-S3 микроконтроллері мен Edge Gateway
арасындағы киберқауіпсіздік мәселелеріне ерекше назар аударылды және
барлық ақпараттық транзакциялар TLS 1.3 хаттамасын пайдаланып шифрланды,
ал аппараттық үдеткіштерді пайдалану жоғары есептеу жылдамдығы мен жалпы
өнімділікті қамтамасыз етті.

{\bfseries 2-кесте. IIoT байланыс хаттамаларының өнімділік көрсеткіштерін
салыстырмалы талдау}

%% \begin{longtable}[]{@{}
%%   >{\raggedright\arraybackslash}p{(\linewidth - 6\tabcolsep) * \real{0.2440}}
%%   >{\raggedright\arraybackslash}p{(\linewidth - 6\tabcolsep) * \real{0.2381}}
%%   >{\raggedright\arraybackslash}p{(\linewidth - 6\tabcolsep) * \real{0.3114}}
%%   >{\raggedright\arraybackslash}p{(\linewidth - 6\tabcolsep) * \real{0.2065}}@{}}
%% \toprule\noalign{}
%% \begin{minipage}[b]{\linewidth}\centering
%% {\bfseries Параметрлер}
%% \end{minipage} & \begin{minipage}[b]{\linewidth}\centering
%% {\bfseries MQTT 5.0 (over TCP)}
%% \end{minipage} & \begin{minipage}[b]{\linewidth}\centering
%% {\bfseries WebSocket (Full-Duplex)}
%% \end{minipage} & \begin{minipage}[b]{\linewidth}\centering
%% {\bfseries HTTP/2 (REST)}
%% \end{minipage} \\
%% \midrule\noalign{}
%% \endhead
%% \bottomrule\noalign{}
%% \endlastfoot
%% Орташа кідіріс (Latency) & 25,4 мс & 38 мс & 72 мс \\
%% Желілік шығын (Overhead) & 12 байт & 8 байт & 40 байт \\
%% Тұрақтылық (Jitter) & Жоғары (±50.2 мс) & Төмен (±25.5 мс) & Өте жоғары
%% (±75.8 мс) \\
%% Қолданылуы & Сенсорлық деректерді жинау & 3D Dashboard/Визуализация &
%% Мерзімді синхрондау \\
%% \end{longtable}

\fig{i2/image23}[]

{\bfseries 7-сурет. IIoT хаттамаларының (MQTT 5.0, WebSocket, HTTP/2)
орташа кідіріс уақыты (Latency) мен желілік артық шығындарын (Overhead)
салыстырмалы талдау}

Digital Twin жүйесінде ақпарат алмасудың тиімділігін бағалау үшін
жүргізілген эксперименттік зерттеулердің нәтижелері 7-суретте
көрсетілген. Графиктерді талдау негізінде келесі қорытындылар жасалды:

а) Кідірісті талдау: эксперименттік зерттеулердің нәтижесінде MQTT 5
протоколы үшін ең төменгі орташа кідіріс белгіленді (25,4 мс). Бұл оның
нақты уақытқа мүмкіндігінше жақын жұмыс істеу қабілетін дәлелдейді.
Сонымен қатар, WebSocket протоколы 38 мс нәтиже көрсетті. Бұл көрсеткіш
оның толық дуплексті сипатына (тұрақты екі жақты байланыс) және қайта
қосылуға кететін уақыттың болмауына байланысты. HTTP / 2 протоколы
сұраныс-жауап моделінің күрделілігіне байланысты ең жоғары кідіріске (72
мс) ие болды.

b) Желілік шығындар: WebSocket деректер пакетінің тақырыбы мен үстеме
шығындар (8 байт) тұрғысынан ең үнемді хаттама ретінде танылды. MQTT 5.0
12 байттық өткізу қабілеттілігін қамтамасыз ету арқылы желілік
ресурстарды пайдалану тиімділігін көрсетті, бұл әсіресе өткізу қабілеті
төмен IIoT желілері үшін маңызды. HTTP / 2 протоколы мәтіндік
тақырыптарды қолдануға байланысты (тіпті қысылған) ең жоғары мәнге ие
(40 байт).

 Талдау нәтижелері WebSocket және MQTT
протоколдары нақты уақыт режимінде "цифрлық егіздерді" 3D
визуализациялау және бақылау үшін ең оңтайлы шешімдер екенін көрсетеді.

 WebSocket ең жоғары жылдамдықты ұсынса, MQTT 5
ең жоғары жылдамдықты қамтамасыз етеді және MQTT 5.0 (QoS 1 деңгейімен)
деректерді шлюзге жеткізуге және трафикті үнемдеуге кепілдік береді.
HTTP / 2 протоколын тек үлкен конфигурация файлдарын мезгіл-мезгіл
синхрондау және тасымалдау үшін пайдалану ұсынылады.

\emph{Өнімділікті бағалаудың математикалық әдістемесі.} Хаттаманың
тиімділігін сандық бағалау үшін келесі математикалық модельдер
қолданылды:

\def\labelenumi{\alph{enumi})}

1. RTT есептеу:

\begin{quote}
\(RTT = \ T_{received} - \ T_{sent}\) (2)
\end{quote}

Мұндағы \(T_{sent}\ \)- микроконтроллер деректі жіберген уақыт
белгісі,\(T_{received\ \ }\)- растауда келетін уақыт.

\def\labelenumi{\alph{enumi})}
\setcounter{enumi}{1}

1. Деректердің өзгерістері:

Желілік тұрақтылықты қалыптасқан ауытқу (\(\sigma\)) арқылы анықтаймыз:

\(\sigma = \ \sqrt{\frac{1}{N}\sum_{i = 1}^{N}\left( x_{i} - \mu \right)^{2}}\ \)
(3)

мұнда \(x_{i}\)- кідіріс мәні, \(\mu\) - орташа мән.

\def\labelenumi{\alph{enumi})}
\setcounter{enumi}{2}

1. Деректерді берудің өткізу қабілеті:

Жүйенің белгілі бір уақыт аралығындағы ақпаратты тасымалдау жылдамдығын
келесі формуламен анықтаймыз:

\(S = \ \frac{D}{T_{total}}\) (4)

мұндағы \(D\) - деректердің жалпы көлемі (бит немесе байт),
\(T_{total}\ \)- деректерді жіберуге кеткен толық уақыт. Бұл көрсеткіш
нақты тиімділікті сипаттайды.

Архитектураның интегралдық тиімділігін бағалау үшін авторлық тиімділік
коэффициенті (\(K_{ff}\)) енгізілді:

\(K_{eff} = \frac{S_{reduction} \bullet (1 - PLR)}{{RTT}_{avg}}\) (5)

Мұндағы \(S_{reduction}\)- деректерді жинақтау деңгейі, \(PLR\) -
пакеттердің жоғалу коэффициенті, \({RTT}_{avg}\) - орташа кідіріс
уақыты. Бұл көрсеткіш жүйенің өткізу қабілеті мен сенімділігінің
арақатынасын сандық түрде сипаттауға мүмкіндік береді.

Орташа RTT: 25,4 мс.

Пакеттердің жоғалуы (PLR): 0,34\%.

Өткізу қабілеті: 125.0 kbps.

Ұсынылған математикалық модель негізінде құрылғының сыни сипаттамалары
есептелді. Эксперименттік деректерді (\(S_{reduction}\) = 67,06\%,
\(PLR\) = 0,34\%, \({RTT}_{avg}\) = 25,4 мс) формулаға орналастыру
арқылы келесі нәтиже алынды:

\(K_{eff} = \frac{67.06 \bullet (1 - 0,0034)}{25,4} \approx 2,63\) (6)

Алынған \(K_{eff} \approx 2,63\) өлшемі таңдалған Protobuf құрылымы мен
MQTT 5.0 протоколы деректер жылдамдығы мен сенімділік арасындағы оңтайлы
тепе-теңдікті қамтамасыз ететінін сандық түрде дәлелдейді.

Пакеттердің жоғалу коэффициенті (Packet Loss Rate):

Деректерді беру кезінде қателер мен жоғалған деректердің үлесін есептеу
үшін төмендегі формула қолданылады:

\(PLR = \ \frac{N_{lost}}{N_{sent}}\  \times 100\%\) (7)

мұнда \(N_{lost}\) - қабылданбаған пакеттер саны, \(N_{sent}\) -
жіберілген пакеттердің жалпы саны.

Жоғарыда келтірілген математикалық модельдер негізінде
микроконтроллердің жұмысын анықтау үшін көптеген сынақтар жүргізілді.
Нәтижелер төмендегі 3-кестеде жинақталған.

{\bfseries 3-кесте. Деректерді беру сапасының көрсеткіштері}

%% \begin{longtable}[]{@{}
%%   >{\raggedright\arraybackslash}p{(\linewidth - 8\tabcolsep) * \real{0.1765}}
%%   >{\centering\arraybackslash}p{(\linewidth - 8\tabcolsep) * \real{0.1365}}
%%   >{\centering\arraybackslash}p{(\linewidth - 8\tabcolsep) * \real{0.1664}}
%%   >{\centering\arraybackslash}p{(\linewidth - 8\tabcolsep) * \real{0.2459}}
%%   >{\centering\arraybackslash}p{(\linewidth - 8\tabcolsep) * \real{0.2748}}@{}}
%% \toprule\noalign{}
%% \begin{minipage}[b]{\linewidth}\centering
%% {\bfseries Өлшеу №}
%% \end{minipage} & \begin{minipage}[b]{\linewidth}\centering
%% {\bfseries RTT (ms)}
%% \end{minipage} & \begin{minipage}[b]{\linewidth}\centering
%% {\bfseries Jitter (σ,ms)}
%% \end{minipage} & \begin{minipage}[b]{\linewidth}\centering
%% {\bfseries Throughput (kbps)}
%% \end{minipage} & \begin{minipage}[b]{\linewidth}\centering
%% {\bfseries Packet Loss (PLR,\%)}
%% \end{minipage} \\
%% \midrule\noalign{}
%% \endhead
%% \bottomrule\noalign{}
%% \endlastfoot
%% 1 & 24 & 1.2 & 128 & 0.0 \\
%% 2 & 26 & 2.1 & 125 & 0.5 \\
%% 3 & 22 & 0.8 & 130 & 0.0 \\
%% 4 & 30 & 4.5 & 115 & 1.2 \\
%% 5 & 25 & 1.5 & 127 & 0.0 \\
%% {\bfseries Орташа мән} & {\bfseries 25.4} & {\bfseries 2.02} & {\bfseries 125.0} &
%% {\bfseries 0.34} \\
%% \end{longtable}

Кестеден көріп отырғаныңыздай, орташа RTT әлі де 25.4 мс деңгейінде. Бұл
нақты уақыт режимінде өңдеу үшін жақсы эталон. Сонымен қатар,
пакеттердің жоғалу коэффициенті 0,3\% құрайды, бұл таңдалған байланыс
протоколының сенімділігі жоғары екенін көрсетеді.

\fig{i2/image24}[]

{\bfseries 8-сурет. RTT және Джиттер көрсеткіштерінің динамикасы}

Көк сызықпен ұсынылған бірінші графикте пакеттің қайтару уақыты немесе
RTT 22-ден 29 миллисекундқа дейін өзгеретіні көрсетілген. Қызыл сызықпен
көрсетілген діріл мәні төртінші өлшемде 4,5 миллисекундта максимумға
жетеді, бұл желіде уақыт кідірістерінің болуын көрсетеді. Бұл өзгеріске
қарамастан, орташа көрсеткіштер жүйенің тұрақты жұмыс істейтінін
байқатады. MQTT 5.0 протоколында байқалған тербелістер оның QoS 1
функциясымен байланысты, бұл әр пакеттің PUBACK хабарламасымен расталуын
қамтамасыз етеді. Бұл растау процесі желі қолжетімсіз болған кезде
пакеттерді қайта жіберуге тырысумен қатар, қалтырау деп аталатын пакет
кірістері арасындағы уақыт айырмашылығын арттырады.

\fig{i2/image25}[]

{\bfseries 9-сурет. Өткізу қабілетінің (Throughput) өзгеруі)}

Бұл гистограмма жазбалардың ауысу жылдамдығын (кбит/с) сипаттайды. Өлшеу
кезінде жылдамдық 118-ден 130 кбит/с-қа дейін өзгерді және сақталды,
жылдамдықтың төмендеуі (4-ші өлшеу) кідіріссіз пакеттің жоғалуына (PLR)
және JIT-тің кең спектріндегі жарылысқа байланысты болды. Математикалық
бағалау және эксперименттік өлшеулер құрылғының жалпы RTT 25,4 мс екенін
растады. Бұл нақты уақыт режимінде жазбаларды толық өңдеуге мүмкіндік
береді. Cонымен қатар, пакетті жоғалту коэффициенті (0,3\%) протоколдың
жоғары сенімділігін, өткізу қабілеттілігі мен дірілдің тұрақтылығын
көрсетеді. Мәндер таңдалған ережелер жиынтығы тиімді жұмыс істейтінін
эксклюзивті шарттарда дәлелдеді.

\emph{Шеткі есептеулер және адаптивті фильтрация алгоритмінің
тиімділігін негіздеу.} Шекті есептеу деңгейінде қауымдастық жүктемесінің
төмендеуі мен бұлттық инфрақұрылымның шамадан тыс жүктелу белгілерін
математикалық және эксперименттік негіздеу үшін шекті мәндерге
негізделген толық жауап беретін сүзу әдісі қолданылды. Бұлтты
инфрақұрылым, шекті мәндерге негізделген толық жауап беретін сүзу
тәсілі, перифериялық есептеу деңгейінде қолданыла бастады. Қауымдастық
жүктемесін 80\% - ға дейін төмендетудің негіздемесі: кәдімгі жүйелерде
сенсор ақпаратты үздіксіз жібереді, мысалы, секундына 1 рет. Бұл
қауымдастықты мағынасыз ақпаратпен (Шу) толтырады. Біз ұсынатын
архитектурада ESP32-S3 микроконтроллері физикалық параметрлер
сипатталған шекті мәннен (θ) асып кетсе, бұлттық серверге мүлдем жаңа
өлшем жібереді.Фильтрация шарты келесідей сипатталады:

\(|\mathrm{\Delta}V| = \left| V_{t} - V_{t - 1} \right| \geq \ \theta\)
(8)

мұндағы \(V_{t}\) - дәл уақыттағы өлшем (температура, ылғалдылық немесе
қысым), \(V_{t - 1}\) - алдыңғы жіберілген өлшем, ал \(\theta\) - рұқсат
етілетін ауытқу шегі (мысалы, температура үшін 0,5\(℃\))

Желілік кептелістің төмендеу коэффициенті \(R_{net}\) келесі
математикалық формуламен есептелді:

\(R_{net} = (1 - \frac{N_{filtered}}{N_{raw}}) \times 100\%\) (9)

Эксперимент қалыпты өндірістік ортада 1 сағат ішінде BME280 сенсорынан
\(N_{raw} = 3600\)өлшемін алды. Адаптивті сүзгілеу алгоритмін қолдана
отырып, шлюзге тек \(N_{filtered} = 720\) пакет (жалпы көлемнің 20\%)
жіберілді. Деректердің қалған 80\% - ы (өзгеріссіз қалған немесе шу
болып табылатын параметрлер) шектен тыс бұғатталған. Осылайша, желінің
шамадан тыс жүктелуі 80\% - ға төмендегені эксперименталды түрде
дәлелденді.

Бұлттық инфрақұрылымға салмақ жоғалтудың 40\% негіздемесі: бұлттық
серверге салмақ тек кіріс трафикпен ғана емес, сонымен қатар дерекқорға
жазу транзакцияларымен (Database) және процессор уақытымен (CPU time)
өлшенеді. Mosquitto брокері ретінде қызмет ететін аталған сервер
деректер ағынын оңтайландырып қана қоймайды, сонымен қатар перифериялық
есептеу деңгейінде ақпаратты алдын ала сүзу функциясын орындайды. Бұлтты
сервердің процессоры мен жедел жадының ресурстарын тұтынуды бақылау
мәліметтер базасына сұраныстар пакеттер санын азайту және Protobuf
форматындағы деректерді капсулалау арқылы жеңілдетілгенін көрсетті.

\(L_{reduction} = \frac{{CPL}_{raw} - {CPL}_{filtered}}{{CPL}_{raw}} \times 100\% \approx 40\%\ \ \ \ \ \ \ \ \ \ \ \ \ \ \ \ \ \ \ \ \ \ \ \ \ \ \ \ \ \ \ \ \ \ \ \ \ \ \ \ \)
(10)

мұндағы \(CPL\) (Cloud processing Load) бұлттық инфрақұрылымда
транзакцияларды өңдеу және сақтау үшін пайдаланылатын өңдеу қуатының
орташа көрсеткіші (CPU/RAM) болып табылады. Бұл бұлтты инфрақұрылымда
транзакцияларды өңдеу және сақтау үшін қолданылады. Бұл көрсеткішті
сүзгісіз режиммен салыстырғанда 40\% - ға төмендету жүйенің
масштабталуын айтарлықтай арттырады.

{\bfseries Нәтижелер мен талқылау.} Жүргізілген эксперименттік сынақтар
4-индустрия контекстінде цифрлық егіздер туралы ақпарат алмасуды
оңтайландырудың нақты тетіктерін анықтады. Зерттеу Protobuf екілік
серияландыру форматының тиімділігін растады: деректер пакетінің көлемі
85-тен 28 байтқа дейін қысқарып, желілік трафикті 67,06\%-ға үнемдеуге
және микроконтроллердің қуат тұтынуын 30-40\%-ға төмендетуге мүмкіндік
берді. Хаттамаларды саралау нәтижесінде MQTT 5.0 датчиктерден деректерді
жинау үшін ең қолайлы екенін көрсетті, ол минималды кідіріспен (25,4 мс)
және желідегі төмен шығындармен (12 байт) сипатталады. WebSocket
протоколы (38 мс кідіріс) нақты уақыттағы 3D визуализациясы үшін тұрақты
арнаны қамтамасыз етсе, HTTP/2 (72 мс) тек файлдарды мезгіл-мезгіл
синхрондау үшін тиімді екенін дәлелдеді. Сонымен қатар, перифериялық
есептеу (Edge Computing) деңгейінде адаптивті шекті сүзуді енгізу
желінің шамадан тыс жүктелуін 80\%-ға және бұлттық инфрақұрылым
жүктемесін 40\%-ға төмендетіп, шлюзге тек маңызды өзгерістерді жіберді.
TLS 1.3 стандартымен қорғалған бұл архитектураның интегралды тиімділік
коэффициенті \(K_{eff} \approx 2,63\) құрап, жылдамдық, қауіпсіздік және
ресурстарды үнемдеу арасындағы оңтайлы тепе

-теңдікті ғылыми түрде растады.

{\bfseries Қорытынды.} Бұл шолудың нәтижелері ресурстармен шектелген IIoT
гаджеттер мен виртуалды егіздер арасындағы ақпарат алмасуын оңтайландыру
мақсатына толығымен қол жеткізілгенін көрсетеді. Тәжірибенің бір
кезеңінде ұсынылған аппараттық және бағдарламалық жасақтама пакеті
(ESP32-S3, BME280, MQTT 5.0 және Protobuf) құрылғының өткізу
қабілеттілігі мен кейінге қалдырылған уақыт арасында жоғары тұрақтылықты
қамтамасыз ететінін растады.

Бейімделгіш сүзгілеу мен екілік сериялау форматтарын араластырудың
арқасында бұлттық инфрақұрылымдағы жүктемені жеңілдету және қоғамдық
ресурстарды үнемдеу міндеттері дұрыс шешілді. Сонымен қатар, TLS 1.3
танымал бағдарламалық жасақтамасы өндірістік штамдардың киберқауіпсіздік
талаптарына толығымен сәйкес келетінін растады. Алынған клиникалық
нәтижелер ақылды кәсіпорындарда нақты уақыт режиміндегі бақылау
құрылымдарын жобалау үшін берік негіз болады. Болашақ зерттеулер бұл
құрылымды ірі масштабты сенсорлық желілерге (торлы желілерге) бейімдеу
және перифериялық құрылғыларда синтетикалық интеллект (TinyML)
алгоритмдерін пайдалану туралы хабардарлықты арттырады.

\emph{{\bfseries Қаржыландыру.}} Мақала Қазақстан Республикасының Ғылым
және жоғары білім министрлігінің қолдауымен, BR24992975 "Жасанды
интеллект және IIoT технологияларын қолдана отырып, тамақ өңдеу
кәсіпорнының цифрлық егізін жасау" бағдарламалық-мақсатты қаржыландыру
жобасы шеңберінде жарияланды, 2024-2026 жж.

{\bfseries Әдебиеттер}

1. Tao F., Zhang, H., Nee A.Y. Digital Twin in Industry:
State-of-the-Art.// IEEE Transactions on Industrial Informatics.- 2019.
- Vol.15(4). - P.2405--2415.
\href{https://doi.org/10.1109/TII.2018.2873186}{DOI
10.1109/TII.2018.2873186}.

2. Postel J. Transmission Control Protocol.// IETF RFC.-
1981. \href{https://www.google.com/search?q=https://doi.org/10.17487/RFC0793}{DOI
10.17487/RFC0793}.
https://dl.acm.org/doi/10.17487/RFC0793\#issue-downloads

3. Fette I., Melnikov A. The WebSocket Protocol.// IETF RFC.- 2011.- 71
p. ISSN 2070-1721.

4. OASIS Open. MQTT Version 5.0. OASIS Standard.// OASIS.- 2019.
\ul{\url{https://docs.oasis-open.org/mqtt/mqtt/v5.0/os/mqtt-v5.0-os.html}.
-} Date accessed: 27.12.2025 .

5. Mishra B., Kertesz A. The Use of MQTT in M2M and IoT Systems: A
Survey.// IEEE Access.- 2020. - Vol.8:201071-201086.
\href{https://doi.org/10.1109/ACCESS.2020.3035849}{DOI
10.1109/ACCESS.2020.3035849}.

6. Rescorla E. The Transport Layer Security (TLS) Protocol Version
1.3. // IETF RFC.- 2018. \href{https://doi.org/10.17487/RFC8446}{DOI
10.17487/RFC8446}.

7. Silva D. R., Oliveira G. M., Silva I., Ferrari P., Sisinni E. Latency
evaluation for MQTT and WebSocket Protocols: an Industry 4.0
perspective.// IEEE ISCC.- 2018. - P.667-1238.
\href{https://doi.org/10.1109/ISCC.2018.8538692}{DOI
10.1109/ISCC.2018.8538692}.

8. Lakshminarayana S., Praseed A., Thilagam P. S. Securing the IoT
Application Layer From an MQTT Protocol Perspective: Challenges and
Research Prospects. // IEEE Communications Surveys \& Tutorials.- 2024.-
P.99.
\href{https://www.google.com/search?q=https://doi.org/10.1109/COMST.2024.3372630}{DOI
10.1109/COMST.2024.3372630}.

9. Cho Y., Noh S. D. Design and Implementation of Digital Twin Factory
Synchronized in Real-Time Using MQTT. // Machines.- 2024. -- Vol.
12(11):759. \href{https://doi.org/10.3390/machines12110759}{DOI
10.3390/machines12110759}.

10. Fei T., et al. makeTwin: A reference architecture for digital twin
software platform.// Chinese Journal of Aeronautics.- 2024. -- Vol.
37(1):1-18. \href{https://doi.org/10.1016/j.cja.2023.09.005}{DOI
10.1016/j.cja.2023.05.002\ul{.}}

11. Khan S., et al. A Novel Digital Twin (DT) Model Based on WiFi CSI,
Signal Processing and Machine Learning for Patient Respiration
Monitoring and Decision-Support.// IEEE Access.- 2023. -- Vol.11.-
P.103554-103568.
\href{https://www.google.com/search?q=https://doi.org/10.1109/ACCESS.2023.3316508}{DOI
10.1109/ACCESS.2023.3316508}.

12. Gao Y., et al. Digital Twin-Based Office Equipment Management and
Personnel Detection System.// IEEE CSCWD.- 2024.- P.1651-1656. DOI
\href{https://doi.org/10.1109/CSCWD61410.2024.10580816}{10.1109/CSCWD61410.2024.10580816}.

13. Skvorc M., Horvat M., Srbljic S. Performance evaluation of Websocket
protocol for implementation of full-duplex web streams.// IEEE MIPRO.-
2014. - P.1003-1008.
\href{https://doi.org/10.1109/MIPRO.2014.6859715}{DOI
10.1109/MIPRO.2014.6859715}.

14. Mijovic S., et al. Comparing application layer protocols for the
Internet of Things via experimentation.// IEEE RTSI.- 2016.
\href{https://www.google.com/search?q=https://doi.org/10.1109/RTSI.2016.7740582}{DOI
\href{https://doi.org/10.1109/RTSI.2016.7740559}{10.1109/RTSI.2016.7740559}.}

15. Nair S. S. Securing Against Advanced Cyber Threats: A Comprehensive
Guide to Phishing, XSS, and SQL Injection Defense. // JCSTS. - 2024. -
Vol.6(1):76-93. \href{https://doi.org/10.32996/jcsts.2024.6.1.9}{DOI
10.32996/jcsts.2024.6.1.9}.

\emph{{\bfseries Авторлар туралы мәлімет}}

Амирханова Г.А.- PhD, әл-Фараби атындағы Қазақ ұлттық университеті,
Алматы, Қазақстан, е-mail:
gulshat.aa@gmail.com;

Фарузқызы Е.- бакалавр, әл-Фараби атындағы Қазақ ұлттық университеті,
Алматы, Қазақстан, е-mail:
elnura.dauletkhan@gmail.com;

Амирханов Б. - жасанды интеллект және Big Data кафедрасы, әл-Фараби
атындағы Қазақ ұлттық университеті, Алматы, Қазақстан, е-mail:
amirkhanov.b@gmail.com;

Раева А.А.- бакалавр, әл-Фараби атындағы Қазақ ұлттық университеті,
Алматы, Қазақстан, е-mail:
alinaraevali18@gmail.com;

Жайсанова Д.С.- әл-Фараби атындағы Қазақ ұлттық университеті, Алматы,
Қазақстан, е-mail:
zhaisanova15@gmail.com.

\emph{{\bfseries Information about the authors}}

Amirkhanova G. - PhD, Al-Farabi Kazakh National University, Almaty,
Kazakhstan, e-mail:
gulshat.aa@gmail.com;

Faruzkyzy E. - bachelor, Al-Farabi Kazakh National University, Almaty,
Kazakhstan, e-mail:
elnura.dauletkhan@gmail.com;

Amirkhanov B. - department of Artificial Intelligence and Big Data,
Al-Farabi Kazakh National University, Almaty, Kazakhstan, e-mail:
amirkhanov.b@gmail.com;

Rayeva A.- bachelor, Al-Farabi Kazakh National University, Almaty,
Kazakhstan, e-mail:
alinaraevali18@gmail.com;

Zhaissanova D. - Al-Farabi Kazakh National University, Almaty,
Kazakhstan, e-mail:
zhaisanova15@gmail.com.